\chapter[Christianity: A World Religion]{Christianity: How a Local Movement Became a World Religion}
\section{An Oil Lamp Engraved with a Fish}
First, pick up an object: a clay oil lamp.

It is small enough to fit in one hand, with a narrow filling hole in its body and a soot-blackened spout: it really was lit, many times. Many such lamps have been excavated from Roman catacombs. Rows of niches cut into the passage walls held the dead; the living carried lamps like these down to visit their graves. What makes this lamp unusual is a little design engraved underneath: a fish.

It is not a realistic fish, just an outline made of a few curves that a child could draw. Why engrave a fish on an everyday object? Because the Greek word for fish, ichthys, consists of five Greek letters that can serve as the initials of a phrase: ``Jesus Christ, God's Son, Saviour.'' To an outsider it is only a fish; to an insider it is an entire declaration of faith. Like a code recognised only by fellow members, it tells the owner of an unfamiliar courtyard, when you arrive carrying your lamp, whether to pour you wine or shut the door.

Remember this design. This chapter begins with the declaration hidden inside that fish: a man executed as a criminal by an empire, followers afraid to identify themselves publicly, a little movement still in the process of writing its own scriptures. How did it become the official religion of the empire that executed him three centuries later, and, two thousand years later, the collective name for churches with over two billion members worldwide, roughly a third of humanity?

What happened in between is far more complicated, and more interesting, than ``a good man moved the world.''

\section{An Evening in Corinth}
Now come with me into a scene.

Time: an evening around 56 CE. Place: Corinth, a port city on the Greek peninsula, where cargo from the gulfs to its east and west is unloaded and transshipped. Sailors, merchants, swindlers, and regional dialects mingle endlessly on the waterfront. It is a rich, noisy city where nobody defers to anybody.

We turn down an alley near the docks and enter a household courtyard. Several oil lamps burn indoors; perhaps one has that fish engraved underneath. A dozen or so people sit on benches and floor mats: a small dye-shop owner, his enslaved female household helper, a merchant from the synagogue, and an illiterate dockworker. By the rules of the Roman world, these people should not be eating at the same table. Enslaved and free people share a table, men and women dine together, Jews and Gentiles drink from one cup: all of it is happening tonight.

They are eating dinner, but it is more than a meal. Someone breaks a loaf for everyone to share; they pass around a cup of wine, and someone quietly says, ``In remembrance of him.'' The custom comes from a supper outside Jerusalem over twenty years earlier, when their teacher was alive.

Afterwards, someone produces a papyrus roll: a letter. Everyone falls silent. Its author, Paul, is a Jewish man who travels all over the empire. Some years ago he spent eighteen months in this city, personally bringing these people into the group. Now elsewhere, he has heard about the Corinthian church's troubles: factions at gatherings, displays of wealth, incest, and furious arguments about eating meat sacrificed to idols. This letter is his reply.

The reader proceeds word by word. There are rebukes, answers, and a passage that will travel around the world: without love, even magnificent eloquence is only a clanging gong; love is patient and kind. At an important passage, the enslaved woman in the corner straightens up. Perhaps she alone in this room has had nobody speak to her like this all day.

Remember this room, its mixed company, the letter travelling hundreds of kilometres and being copied along the way, and the supper held ``in remembrance of him.'' The rest of this chapter asks how a handful of rooms like this along the eastern Mediterranean became thousands, then millions.

\section[Why Did This Failed Movement Survive?]{The Real Question: Why Did a Loser's Movement Not Disappear?}
First fill in the background, then state the real puzzle. Do not rush to an answer.

The movement's central figure was Jesus, a Jew. This must come first because later stories so easily obscure it: Jesus was born in Jewish lands, read Jewish scriptures, entered Jewish synagogues, and observed Jewish festivals. His opponents, scriptural quotations, and parables all belonged to Jewish tradition. Jews then lived under Roman rule and within an ancient hope: they believed God had promised their ancestors to send an ``anointed one,'' called the ``Messiah'' in Hebrew and ``Christ'' in Greek, to restore their people. A small group of disciples gathered around Jesus and declared that he was this person. The claim alarmed Jerusalem's religious leaders and the Roman governor. You know the outcome: crucifixion, a Roman punishment reserved for enslaved people, rebels, and serious criminals, deliberately public, prolonged, and humiliating, to make an example of its victim.

The disciples scattered to save themselves. So far, this story looks no different from countless failed uprisings. Jesus was far from the only ``prophet'' or ``saviour'' the Romans executed; those movements all disappeared with their leaders' deaths.

Normally this one should have done so too. Instead, something strange happened: a few weeks later, the scattered disciples reassembled in Jerusalem and publicly claimed that Jesus had risen from the dead. Christians believe this was a real miracle, the foundation of their faith. Historians can establish only something else: these people genuinely believed it, strongly enough to die for it. Throughout this chapter, keep these two statements separate: ``what believers believe'' and ``what historical evidence can show.''

That was not the only surprise. This began as an argument within Judaism: a failed candidate for Messiah, not fundamentally different from thousands of synagogue disputes. Yet within decades the internal argument moved beyond the Jewish people and Palestine, spreading along Roman roads and shipping routes throughout the Mediterranean. Jews passed it to Gentiles: a members-only restaurant suddenly announced free admission for the world. Some crucial rewriting must have happened in between.

Our real question now emerges as a chain of questions:

How did a movement ending in a public execution turn its fortunes around? How did an internal Jewish dispute become a religion open to everyone? How were the particular hopes of one particular people translated into language Greeks, Romans, and Egyptians could make their own? And when we say ``Christianity,'' do we mean beliefs, organisations, ordinary people's lamps and festivals, or emperors' edicts? Are these really the same thing?

The last question is the least conspicuous but the most important. We will devote an entire section to it.

\section[Governors, Synagogues, and Outsiders]{Many Voices: The Governor, the Synagogue, Gentiles, and an Enslaved Woman}
The most honest approach to a historical episode is to listen to several contradictory accounts at once. At least four voices spoke about this emerging group.

\textbf{First voice: a Roman official's report.} Around 112 CE, Pliny the Younger, a Roman aristocrat governing a region in present-day Turkey, wrote to Emperor Trajan for instructions. The letter survives. In essence: Christians have been arrested in my province; I have never handled such cases and am unsure of the rules. Should all be executed, or only those guilty of other crimes? I question them three times, execute those who persist in acknowledging their faith, and release those willing to sacrifice to our images and curse Christ. He added that this ``superstition'' had spread through cities and villages, even depressing business at the temples. The emperor's surviving reply says, in essence: do not hunt them down; act only against those accused who refuse to recant; accept no anonymous accusations.

Pause for a serious exercise: \textbf{state the Roman case in full.} Textbooks often portray persecuting Romans as outright villains. From Pliny's position, however, things looked like this. Rome did not separate religion and government; nobody around the Mediterranean then imagined such a separation. Sacrificing to the emperor functioned rather like a modern civic oath: it tested recognition of the community, not inward belief. These people refused. Worshipping one god and rejecting all others earned them the insult ``atheists'': notice that the term then targeted Christians, not pagans. They also met at night behind closed doors to ``share flesh and blood,'' as outsiders' rumours put it. They would not inform, compromise, or recant even under sentence of death. A responsible official saw a secret association spanning the empire, refusing loyalty ceremonies, unmoved by punishment. You need not think him right---burning any believer today is an atrocity---but you should recognise the intelligible core of his reasoning: he feared not doctrine but loyalty acquiring a second address. Understanding how the machine worked does not excuse what it did to those it crushed.

\textbf{Second voice: neighbours in the synagogue.} To Jerusalem's Jews, the earliest Jesus-followers were not ``people of another religion'' but ``fanatical relatives in our synagogue.'' Both sides read the same scriptures and disputed their interpretation. In 70 CE, Roman troops crushed a Jewish revolt and destroyed Jerusalem's Temple, the centre of the Jewish world. Judaism itself then reorganised amid the ruins, and boundaries between congregations slowly hardened. Some Jewish believers continued observing the Sabbath and attending synagogue while believing Jesus was the Messiah; increasingly, others thought those commitments could no longer fit through one doorway. Separation was not announced on a single day: relations cooled gradually over decades. Remember that our easy modern boxes, ``Judaism'' and ``Christianity,'' did not exist in the first century. There was a roomful of arguing relatives.

\textbf{Third voice: a Gentile mocker.} In the second century, a Greek intellectual called Celsus attacked Christianity in writing. His original is lost, but the Christian scholar Origen quoted extensively while answering him point by point, allowing us to reconstruct it. In essence: these people dare not debate publicly; they target enslaved people, women, children, and the uneducated poor, luring them with ``even fools can be saved.'' His contempt inadvertently preserved important testimony for historians: many early members did come from society's lower ranks and margins. What the mocker most wanted to trample deserves the researcher's closest attention.

\textbf{Fourth voice: arguments inside the movement.} Gentile arrivals immediately sparked uproar. Must a Greek wishing to follow Jesus first accept circumcision and Jewish dietary law? The Jerusalem elders associated with Peter and James argued fiercely with the globe-trotting Paul. Again, practise stating the other side's case fully. Jerusalem's position was not simply stubborn conservatism: the Law embodied centuries of covenant between their people and God; discarding it meant betraying their ancestors and identity. Paul's argument was equally strong: if admission required becoming Jewish first, the movement would remain a Jewish faction and talk of ``all nations'' would be empty. Acts records a Jerusalem council addressing this dispute. Broadly, Paul's approach prevailed: Gentiles need not become Jews first. This internal ruling opened the first gate towards a world religion, detaching faith from ethnic identity and giving it wheels so it could travel.

Together these voices make a striking picture. Outsiders saw a dangerous, strange group; insiders saw outsiders as blind and violent; within the group people still argued furiously over who ``we'' were. History is never ``Christianity'' standing alone in a spotlight and speaking. It is this roomful of noise.

\section[Thinking Bridge: Four Strands]{Thinking Bridge: Unravelling a Religion into Four Strands}
Now introduce our tool. First consider a mistake you have probably made: people discuss religion as one solid object. ``Christianity is peaceful'' or ``Christianity started wars''; ``Buddhism renounces the world'' or ``Buddhist temples owned vast estates.'' These categorical statements all make the same mistake: twisting four different things into one rope.

\textbf{First strand: doctrine.} A religion's officially declared central beliefs. For Christianity these concern God, Christ, redemption, and resurrection. Written in creeds and theology, doctrine answers, ``What does this religion believe?''

\textbf{Second strand: institutions.} Who has authority and how organisation works. How are bishops appointed? Does the pope have supreme authority? How does the church collect taxes, run schools, or vote at meetings? Institutions answer, ``How does this religion operate?''

\textbf{Third strand: popular practice.} What ordinary believers actually do, whatever the books say: lighting candles, hanging amulets, celebrating festivals, engraving fish on lamps. Popular practice answers, ``What does this religion look like in everyday life?''

\textbf{Fourth strand: political history.} What happens between religious organisations and governments, wars, and economies. Why do emperors support them? Whom have armies fought under their banners? Political history answers, ``What worldly affairs has this religion become involved in?''

The strands intertwine but are not identical. Within the same ``Christianity,'' doctrine debates God's nature, bishops struggle for institutional power, a peasant woman ties red cloth to a saint's image in popular practice, and an emperor uses the church as political glue for his empire. More importantly, \textbf{events on one strand often cannot settle questions on another.} Medieval armies campaigning in Christ's name, on the political strand, prove neither the falsity nor the truth of the doctrine ``love your enemies'': those are different matters. A corrupt pope, on the institutional strand, does not make an elderly village woman's prayers, on the popular-practice strand, fraudulent.

Defenders make the same substitution in reverse: confronted with institutional wrongdoing, they shift the subject to lofty doctrine. That too confuses the strands.

The tool works beyond Christianity. Next time an online argument declares that some religion ``is inherently violent'' or ``purifies people's hearts,'' ask which strand the claim concerns and which strand supplies its evidence. Thousands of abusive replies may result simply from opponents standing on different strands and swinging at each other across the gap.

Remember this tool: nearly every major split and controversy ahead involves four strands tangling together while pulling in different directions.

\section[Primary Sources: Coin and Creed]{Reading Primary Sources: From a Coin to a Creed}
Now read some original material. The movement's documents, the twenty-seven writings later called the New Testament, were all composed in the first century and have long been in the public domain. The Chinese source quotes the widely used Chinese Union Version, published in 1919 and likewise in the public domain; those quotations are translated below.

\textbf{First passage: a boundary drawn in one sentence.}

Mark recounts an attempt to trap Jesus by asking whether taxes should be paid to the Roman emperor. Saying yes would offend his people, who hated Rome; saying no would invite immediate denunciation for rebellion. Jesus requested the tax coin, a silver piece bearing the emperor's portrait, and asked whose image and inscription it carried. ``Caesar's,'' they answered. Then came words quoted repeatedly for two thousand years:

\begin{quote}
Jesus said, ``Give Caesar's things to Caesar, and God's things to God.'' (Mark 12:17, translated from the Chinese Union Version.)
\end{quote}
At first this seems a clever escape. Think further: for the first time in human history, the sentence declared this plainly that imperial jurisdiction had limits. The purse is yours; some things are not. In an ancient world where politics and religion were inseparable, it drew a previously nonexistent line on the wall. Later Christians stood on that line when refusing imperial sacrifices; European churches and kings stood on it during their struggles. An answer to a trap became the seed of a whole political tradition. Here is a specimen of doctrine quietly transforming political history a thousand years later.

\textbf{Second passage: a declaration and a maths problem.}

John opens not with a story but with a philosopher's declaration:

\begin{quote}
The Word became flesh and lived among us... (John 1:14, translated from the Chinese Union Version.)
\end{quote}
``The Word became flesh'': Christian tradition calls this the \textbf{Incarnation}. Believers hold that God did not watch humanity from afar but became a particular human being who hungered, cried, and died. Alongside it stands \textbf{redemption}: tradition holds that Jesus' death on the cross was not an ordinary execution but repayment of humanity's unpayable debt, reconciling humanity with God. Whether these claims are true remains disputed between believers and sceptics. That these beliefs produced real historical consequences is not disputed.

Trouble followed. If Jesus is God, are there one or two gods? If he is God's Son, is he equal to the Father or slightly lower? Greek philosophy's precise vocabulary entered the battlefield. Early in the fourth century, an Alexandrian teacher, Arius, argued that the Son was created and inferior to the Father. His opponents replied that this would invalidate redemption through the cross. Street arguments escalated into clashes between bishops across the empire. In 325, Emperor Constantine, who had just legalised Christianity, could no longer tolerate cracks in his imperial glue and summoned bishops to Nicaea, in present-day Turkey, to settle matters. The resulting Nicene Creed declared the Son ``begotten, not made, of one substance with the Father'': a tiny verbal difference marked a sharp boundary. Further elaboration produced the doctrine of the \textbf{Trinity}: believers hold that God is one, existing in three persons, Father, Son, and Holy Spirit. Whether this was mystery or contradiction was debated even among Christians. Notice its historical function, however: one council and one document drew the first hard, empire-wide boundary around ``who is a Christian.'' Doctrine, institutions (the bishops' council), and politics (imperial summons and pressure) tied a knot at this moment. Our four-strand tool makes that knot clearly visible.

\textbf{Third passage: admission for those at the bottom.}

Return to the room in Corinth. Why could enslaved people and women feel they belonged? In a letter, Paul wrote words shocking to his contemporaries:

\begin{quote}
There is no distinction between Jew and Greek, free and enslaved, male and female, for you have all become one in Christ Jesus. (Galatians 3:28, translated from the Chinese Union Version; its term ``Hellenes'' means Greeks.)
\end{quote}
Do not mistake this, through modern eyes, for ordinary inspirational comfort. Ethnicity, free status, and gender formed three impregnable walls determining everything in that world. This sentence abolished all three inside the community. It did not immediately free a single enslaved person: institutional equality would take over another millennium, arriving with the church's bloodstains and evasions still attached. But as a community's account of itself, it made admission extraordinarily cheap. You need not be Jewish, male, free, or literate. Compared with Mediterranean religious clubs restricted by ethnicity or status, it dismantled an entrance barrier. Celsus's contempt for ``recruiting the lower classes'' and believers' pride in ``all becoming one'' described opposite sides of the same fact.

\section{Why Did It Prevail? Three Explanations}
The broad facts are uncontroversial. From a few dozen people in the first century, Christians became a substantial share of the empire's population by the early fourth: historians commonly suggest anything from a few per cent to around a tenth, although exact figures are unrecoverable. By the late fourth century, Christianity was the official religion; today it has over two billion adherents worldwide. \textbf{Why?} Here are three genuinely different explanations, each with strong evidence and limits.

\textbf{Explanation one: the community advertised itself (sociology).}

These scholars examine rooms like Corinth's. Early Christian groups offered things as scarce as oases in contemporary cities: support for widows, homes for orphans, and, during epidemics when urban residents fled, Christians who stayed to bring patients water and food. Nursing improved survival and impressed onlookers. Add the absence of entrance barriers and a Mediterranean-wide network connected by letters and hospitality, and the community becomes an efficient machine for gaining members: newcomers received real care, stayed, and brought others. This explanation is concrete, every link rooted in everyday logic. Its limitation is that faith begins to resemble a mutual-aid society. Warm communities alone do not become world religions; many exist, so why did this one expand? It explains attachment better than inspiration.

\textbf{Explanation two: the emperor decided (politics).}

Others focus on the summit of power. Constantine reportedly saw a vision before the Battle of the Milvian Bridge in 312; the following year he issued the Edict of Milan, giving Christianity legal status. In 380, Theodosius I officially made it the imperial religion. Thereafter non-Christians could not hold high office, pagan temples closed, and resources poured into churches. This explanation is blunt: administrative imperial force, not angels, produced world-religion status. But consider a counterexample easily overlooked: \textbf{Armenia made Christianity its state religion around 301, over a decade before Constantine's conversion and beyond Rome's direct control.} Ethiopia's ancient church likewise did not grow under Roman protection. Equating Christianity with an official Roman product is therefore at least a shortcut. Another limitation: Constantine did not choose from a vacuum. He backed a movement already rooted throughout the empire, with organisational discipline surpassing other voluntary associations. The emperor sailed with an existing wind; he did not conjure it into being.

\textbf{Explanation three: the message could travel (intellectual history).}

This approach examines the message's design. Paul's admission reform detached faith from ethnicity; Koine Greek was the eastern Mediterranean's common language; Roman roads and shipping could carry a letter from Syria to Spain within weeks. The resurrection of a defeated man and the promise that everyone could become one addressed ordinary people's spiritual emptiness in the later empire. In a vast, indifferent, unequal world, the message said: you are seen, you are loved, death is not the end. This explains why the message travelled faster than rival religious messages. Its limitation is intellectual historians' temptation to credit everything to ideas, overlooking hands tending plague victims and imperial swords. However good a message, without people practising it through water and food or institutions supporting it, it may remain another fine phrase dissolving into a marketplace.

Each explanation holds part of the truth. Sociology explains growth from below, politics decisions from above, and intellectual history the message's circulation. Treating them as mutually exclusive answers wastes what each offers.

One misleading picture also needs correction: many imagine three uninterrupted centuries of bloody Roman persecution. Evidence instead shows intermittent, regional persecution. Some reigns were bloody; other decades were peaceful. Pliny's uncertainty was more typical of officials. The uninterrupted horror story was intensified by the later church's cultivation of martyrdom's memory. Recognising this does not excuse persecution: even one execution for belief is an atrocity. It makes the question of success more realistic: a group occasionally suppressed but usually half-tolerated had time to weave its network.

\section[Further Challenge: Growth and Division]{Further Challenge: How Small Groups Become Giants and Keep Splitting}
The preceding three sections considered what happened and why. Now bring in a stronger tool from sociology, the study of how society operates. It explains one of Christianity's most striking features: \textbf{why has this religion kept growing while repeatedly dividing?}

In the early twentieth century, Max Weber and Ernst Troeltsch proposed a famous distinction between two ideal types of religious organisation. A \textbf{sect} is small, demanding, and high-tension: individuals choose to join, requirements are strict, relations with the surrounding world are strained, and ``us'' and ``them'' are sharply separated. A \textbf{church} is large, accessible, and low-tension: you are born into it without choosing; accommodating everyone's tastes means compromise, cooperation, and mutual penetration with the world.

Apply this distinction to our story and a recurring mechanism appears.

First-century Jesus-followers were a classic sect: few members, high costs, hard boundaries, a crucified teacher, arrested believers, and stubborn refusal of imperial sacrifice. Tension was maximal. Then came growth, legalisation, and establishment as the state religion. What does establishment mean? Everyone counts as Christian, including unbelievers, the indifferent, and careerists. Barriers and tension fall to zero. The sect becomes a church.

But people's need for high tension does not vanish; it reappears elsewhere. Again and again, voices inside the church protest: it is too comfortable, too worldly, too disgraceful! Real faith should be harder, purer, more exacting. They form small groups, restoring demanding membership and tension: a new sect. It grows, becomes wealthy and comfortable, and splits again. Third-century dissenters from compromise became desert ascetics; medieval monastic orders repeatedly revived ``pure faith''; in sixteenth-century Germany, the monk Martin Luther posted his Ninety-Five Theses against selling indulgences, payments to reduce punishment for sin, setting off the Reformation. Reformers advocated ``justification by faith'': salvation through faith itself, not church-administered rituals and merit. The split produced today's Protestantism, itself an extraordinary accelerator of the cycle, generating hundreds or thousands of denominations: Baptists, Methodists, Quakers, Pentecostals... Each began by promising a return to original purity.

This is neither a moral tale of church corruption nor a theological tale of truth defeating error. It is a structural mechanism: low tension enables scale; scale brings compromise; compromise generates small high-tension groups; those groups are either suppressed or become new large low-tension churches. Doctrinal arguments, such as the indulgence dispute, provide sparks, but the mechanism itself belongs to institutions and sociology.

The mechanism also illuminates another great split, from a different direction. The Roman world's eastern and western halves grew further apart linguistically and culturally: Latin in the west, Greek in the east. In 1054 their churches formally separated in the East--West Schism. The west became \textbf{Catholic}, recognising the pope's supreme authority; the east became \textbf{Eastern Orthodox}, with self-governing churches rejecting papal supremacy. One apparent trigger was the filioque dispute: does the Holy Spirit proceed from the Father ``and the Son,'' or from the Father alone? One Latin word was at stake. Unravel the strands: doctrine did dispute a word; institutions disputed papal decision-making versus collective episcopal deliberation; politics placed a west whose empire had fallen to barbarians opposite a surviving Byzantine Empire. Their churches already lived in different worlds. One word cannot carry a thousand years of grievances; it was only the final straw.

The sect cycle still turns. Since the twentieth century, Christianity's fastest-growing sector has been the \textbf{Pentecostal--charismatic movement}, outside Catholicism, Orthodoxy, and traditional Protestantism. Its intense worship, experiences of miracles, and personal rebirth make membership emotionally demanding and tension extremely high: a textbook sect-type organisation. Where does it grow fastest? Not Europe, with its many empty old churches, but Latin America, sub-Saharan Africa, and Asia. Brazil, Nigeria, and the Philippines now have tens of millions of Pentecostals; Seoul has what is claimed to be the world's largest single congregation. Christianity's global centre of gravity has quietly shifted from Europe and North America to the southern hemisphere.

Finally, consider facts often omitted from Eurocentric accounts, lest you think this religion's history consists of Rome plus western Europe:

\begin{itemize}
\item A stele erected in Xi'an in 781 records that Christianity of the Syriac tradition, then called Jingjiao, the ``Luminous Teaching,'' had reached Tang China along the Silk Road over a century earlier and received imperial favour. This ``Stele of the Propagation in China of the Luminous Religion of Daqin'' survives, predating the Reformation by eight centuries.
\item Kerala, on India's southwest coast, has a Christian community so old that its origins are difficult to establish. Its tradition credits Jesus' disciple Thomas with crossing the sea to found it. Whether that can be verified or not, this community, combining Syriac worship with Indian social customs, has existed for over a thousand years.
\item The Ethiopian Orthodox Church has its own calendar, saints, and churches carved from solid rock: an African Christianity never ``introduced'' by European missionary movements.
\end{itemize}
Together these examples show that Christianity was never simply a parcel dispatched worldwide from Rome. It resembles an open-source system repeatedly rewritten locally, intertwining with languages, festivals, and ancestral memories to produce versions even Rome's bishop might not recognise.

\section[Christmas Eve Apples]{Christmas Eve Apples: The Old Question in Your Classroom}
Return to today, and your school.

On 24 December, your class WeChat group may feature beautifully wrapped apples as gifts. ``Apples on Christmas Eve bring peace'': ping in pingguo, ``apple,'' sounds like ping in ping'an, ``peace,'' and Chinese calls Christmas Eve ``Peaceful Night.'' Notice this little custom. It has no biblical basis; no European or American church has heard of it, and Christians outside China know nothing of it. It grew spontaneously on China's popular-practice strand, plugging a Chinese pun into the framework of an imported festival.

Open a shopping app: Christmas trees, Santa hats, ``Christmas sales.'' Browse social media and the annual argument appears: ``Should Chinese people celebrate Christmas?'' One side calls foreign festivals a betrayal of one's roots; the other asks why a bit of festive fun needs ideological scrutiny.

Unravel the dispute with our tool. Most celebrants stand on the \textbf{popular-practice strand}: lights, gifts, meals, and an occasion in late December, with little connection to doctrine. Opponents play the cards of \textbf{political history} and national identity. Churchgoers celebrate through doctrine and institutions. Three sides fight from different strands without persuading one another because they are not disputing the same thing. A Roman believer refusing imperial sacrifice and a school pupil posting Christmas apples two thousand years later are bundled under one word, ``Christmas'': that is the power of confusing strands.

A deeper echo remains. The sect--church mechanism runs at full speed online in another arena: devoted niche fans versus mainstream popularity, an underground band versus its supposedly corrupted television-era self, a subculture's ``original spirit'' versus commercialisation. ``You have changed; you are no longer pure'' and ``we cannot survive without reaching more people'' replay an ancient contradiction ten thousand times. Next time a community erupts over losing its character, remember Corinth's lamplit room, the council chamber at Nicaea, and Luther's sheet of paper. Humanity has been wrestling with the same problem of organising its ideals for two thousand years.

\section{Your Question: Where Draw the Boundary?}
Return to the little fish engraved on the lamp.

Its holder lived in a time when people needed codes to recognise one another. The fish meant that the boundary must remain clear even at the cost of death: not one step backwards before the incense of imperial sacrifice. Three centuries later, churches of the same faith stood on the empire's grandest squares, their boundaries blurred. Unbelievers, emperors, money, and power all entered. Christianity thereby changed the world and was changed by it.

Its history is a repeatedly redrawn boundary. Too rigid, and it remains a tiny Palestinian heresy nobody remembers today; too loose, and it becomes imperial wallpaper nobody takes seriously. Paul dismantled the ethnic gate; Luther redrew institutional walls; desert hermits maintained their line through solitude; Pentecostal believers in Latin America's corrugated-roof buildings draw theirs around fervent experience. Each redrawing brings cheers of ``at last, back to the beginning'' and sighs of ``it is no longer what it was.''

So here is your question, about more than religion:

\textbf{Where should a small group sincerely hoping to change the world draw its boundary?} Does preserving purity mean accepting marginality? Does reaching the public inevitably mean dilution? Is there a line that neither locks the group inside the catacombs nor dissolves it into the square's noise?

Give your answer and reasons. First weigh what you already hold: the Roman official's not wholly unreasonable fear, Jerusalem's not merely stubborn loyalty, the hands bringing water and food, the emperor's pen exploiting an existing current, and the sect--church mechanism running for two thousand years.

There is no standard answer. But next time a late-night group chat leaves you red-faced in an argument over whether a community has lost its character, remember: the people in that Corinthian room in 56 CE had already drawn the line you are redrawing.
